\section{The Funge And Crank Gauge}

The native count cannot remain a loose tally. A finite grammar must say where
the tally is carried, which distinctions preserve it, and which distinctions
turn it. This is the work of the thirty-six-gate gauge. A gate is not a
mechanical gadget in the argument. It is a finite door through which a
reading must pass if it is to count as part of the gauge. Each gate asks:
does this reading preserve the native pool, or does it select an asymmetric
turn?

The inherited names are funge and crank. A funge is a pooling move. It
keeps a reading with the native matter count. It preserves, gathers, or
absorbs a distinction without forcing a positive sign. A crank is a
selecting move. It turns the comparison so that an asymmetric sign can be
read. The names matter less than the partition they enforce. Every gate must
be assigned to one side: pooling or selecting, matter-preserving or
asymmetry-turning.

This partition continues the whole-and-parts grammar from the foundation. A
pooling move reads the whole. It says that this gate belongs with the native
matter ledger and does not expose a separate positive sign. A selecting move
reads the part. It exposes the gate at which the symmetric baseline no
longer absorbs the distinction. The two moves are not rival mechanisms. They
are the two obligations of a finite gauge that must both preserve a count and
notice when the count has been turned.

The native electron corner is the all-pooling corner. In that reading, every
gate funges. The gauge finds no crank. No gate selects the positive sign, so
the antimatter count remains zero. This is the finite-gauge form of the
native asymmetry: matter is not merely named as present; the whole gauge has
pooled the reading without finding an asymmetric turn.

The paired tilted reading is one crank away. Thirty-five gates still pool
with the matter reading, and one gate selects the positive sign. The exact
number is important because it prevents the positive sign from becoming a
fog. The pair does not dissolve the whole gauge into antimatter. It turns one
gate. The gauge can therefore count matter and antimatter together without
losing the native baseline: thirty-five pooling moves, one selecting move,
thirty-six gates total.

Schematically the closure reads
\[
  N_{\mathrm{pool}} + N_{\mathrm{select}} = 36.
\]
Again the formula only records the grammar. The point is that no gate is
left unclassified. If a gate pools, it belongs to the matter-preserving side.
If it selects, it belongs to the asymmetric side. A hidden third option would
let the gauge carry an unread distinction. The grammar forbids that. The
gauge closes when pooling and selecting account for the whole finite list.

This closure also protects the antimatter count. In the native reading,
\[
  N_{\mathrm{select}} = 0.
\]
In the one-crank pair,
\[
  N_{\mathrm{select}} = 1.
\]
The positive sign therefore appears as a finite selection, not as an
unbounded mirror world. The grammar can say exactly where the asymmetry has
entered: at the selecting gate. Everywhere else the reading remains pooled
with the native matter ledger.

The gauge should not be mistaken for a mere classification table. A table
can list names without saying what they carry. The gauge partitions moves.
It says which operations preserve the native reading and which operations
turn it. A color-like slot, a flavor-like kind, a phase-like transport, and
a magnitude-like band may all be carried through gates, but each carried
label must answer the same question: does it pool or select for this
reading?

The finite bit-vector image is helpful only in this grammatical sense. A bit
marks whether a gate has pooled or selected. The vector is the whole finite
ledger of such marks. It does not replace the argument with a notation. It
shows why the count is auditible: a reader can walk the finite gauge and ask
of each gate whether the native matter reading has been preserved or the
asymmetric sign has been selected.

The superconducting-pair illustration belongs here as a later physical
shadow. A pair can carry exactly one positron-like turn through the gauge
while the remaining gates continue to pool. The example should not be asked
to prove the partition. It shows why a one-crank reading is meaningful: the
positive sign can be admitted without making the entire finite gauge
positive. One gate turns; the rest preserve the native pool.

This partition also refines the matter/antimatter language. Matter is read
off the pooling side of the gauge. Antimatter is read off the selecting side.
Together they exhaust the gauge only after the partition has assigned every
gate. The grammar therefore avoids two mistakes. It does not treat matter
and antimatter as substances poured into a common box. It also does not let
the native matter reading hide the possibility of a selected positive sign.
Pooling and selecting are both explicit.

The thirty-six-gate gauge is thus the structural home of the native
asymmetry. Relative velocity selected a quotient. Holonomy closed a sign.
The native baseline counted matter and no antimatter. The gauge now
partitions the finite doors through which those readings pass. When all
doors pool, the native matter reading is exact. When one door cranks, the
positive sign is selected at a finite place.

The chapter's last turn follows naturally. Once every gate asks whether it
pools or selects, the grammar has a finite family of predicates. Those
predicates cannot remain abstract. They turn back on the reader. The final
station is the reader asked to answer the questions the gauge has exposed.
